\section{Patient and context factors}\label{sec:patients}
\paragraph{Comorbidity, chronicity and adversity} Comorbid personality disorder more than doubles the odds of a poor depression outcome (OR 2.16), whatever the treatment modality \citep{newtonhowes2014personality}; childhood maltreatment predicts non-response or non-remission (OR 1.43; \citealp{nanni2012maltreat}); and psychotherapy for chronic depression has a small effect against controls ($d=0.23$) and did worse than medication in direct comparisons ($d=-0.31$, attributed to dysthymia studies; \citealp{Cuijpers2010Chronic}). These are prognostic factors for all treatments, not demonstrations that CBT specifically fails. For anxiety, overall diagnostic comorbidity was generally unrelated to effect size across 148 treatment samples, with diagnosis-specific exceptions \citep{olatunji2010comorbid}. Trauma-focused CBT for PTSD in military samples shows large within-group effects ($d$ 0.78--1.10) yet about two thirds of patients retained the diagnosis, and the programmes were only marginally better than non-trauma-focused active controls \citep{steenkamp2015military}. Transdiagnostic CBT is one response to comorbidity \citep{newby2015transdiagnostic}. Baseline severity does not moderate CBT versus medication \citep{Weitz2015Severity}, and CBT still helps medication-refractory depression \citep{Wiles2016CoBalTLong}.

\paragraph{Neurodivergence, cognitive impairment, age and medical illness} The evidence on autism is thin rather than negative: ten randomized trials in autistic adults (n = 537) leave most effects very uncertain \citep{perezcastilla2026autistic}, and an earlier mixed-age meta-analysis found that effects depended on who rated them ($g=0.24$ self-report, non-significant; 0.66 informant; 0.73 clinician; \citealp{weston2016autism}). In intellectual disability, pooled effects were small ($g=0.43$, and 0.36 in randomized trials; \citealp{tapp2023id}); in dementia, CBT-family therapies had a small effect on depression (SMD $-0.23$) and an uncertain one on anxiety \citep{orgeta2022dementia}. Age itself is not a barrier: older adults recovered more often in routine IAPT data (OR 1.33; \citealp{saunders2021older}), and late-life anxiety CBT remitted 54\% of patients, though about a third relapsed \citep{johnco2026olderanxiety}. Medical comorbidity attenuates effects: CBT for chronic pain gives very small to small benefits (SMD $-0.09$ to $-0.34$; \citealp{williams2020chronicpain}), and insomnia CBT works but is weaker with medical than psychiatric comorbidity \citep{wu2015cbtimedical}. In substance use, CBT beats minimal treatment but not other specific therapies \citep{magill2019cbtsud}. We found no meta-analysis of educational attainment or IQ as a moderator. Whether adult ADHD moderates CBT outcome is addressed in Section~\ref{sec:moderators}.

\paragraph{Culture, representation and access} Culturally adapted CBT outperforms other conditions ($g=0.67$) and unadapted versions ($g=0.52$; \citealp{hall2016culturalmeta}; $d=0.32$ in \citealp{benish2011myth}, where adapting the explanation of the illness was the one significant moderator); depression trials show similar gains \citep{anik2021culturaldepression}, and CBT transports to Chinese-descent samples with modification \citep{li2023chinese}. But the trial base is narrow: across 342 depression trials, linguistic minorities were far more likely to be excluded than included and effects were almost never reported by ethnic group \citep{polo2019diversity}, in line with the broader WEIRD-sampling critique \citep{henrich2010weird}. Access shapes outcomes: in IAPT, waiting time and missed appointments predicted worse outcomes and service organisation explained 33\% of the variance in reliable improvement \citep{Clark2018IAPT}; recovery was lower in poorer areas \citep{delgadillo2016poverty}. In low-resource settings, lay-delivered CBT-family therapy works (adjusted mean difference $-4.86$ on the SSQ-14 in Zimbabwe, \citealp{chibanda2016friendship}; $-7.57$ on the BDI-II in India, \citealp{patel2017hap}; OR 0.22 for major depression in Pakistan, \citealp{rahman2008pakistan}), so there CBT fails through reach, not efficacy, with the pooled evidence rated mostly low certainty \citep{vanginneken2021primary,barbui2020lmicumbrella}.

\paragraph{Rumination and worry} Rumination predicts the onset and course of depression and overlaps with worry \citep{nolenhoeksema2008rumination}. They are also targets: rumination-focused CBT improved residual depression in a small trial with mediation by rumination change \citep{watkins2011rfcbt}, and repetitive-negative-thinking training reduced 12-month depression and anxiety diagnoses against a waitlist \citep{topper2017worryrumination}. This suggests a process that standard restructuring may miss, not a demonstrated failure.

\paragraph{Reading} R7 and R8 are supported in a limited sense: patient and context factors shift outcomes, but nearly all are shared with other treatments, and the strongest CBT-specific gaps are evidence gaps (autism, cognitive impairment, minority groups) rather than measured failures.
